% !TEX program = pdflatex
% !TEX root = teoria_densidade.tex
% !TEX spellcheck = pt_BR
%%%%%%%%%%%%%%%%%%%%%%%%%%%%%%%%%%%%%%%%%%%%%%%%%%%%%%%
%%   Relatorio Tecnico - Power Density Tool           %%
%%   (usa a MESMA classe/preambulo do template UFC)   %%
%%%%%%%%%%%%%%%%%%%%%%%%%%%%%%%%%%%%%%%%%%%%%%%%%%%%%%%
\documentclass[
	a4paper, 12pt,
	chapter=TITLE, section=Title, subsection=Title,
	oneside, english, brazil, fleqn
]{abntex2}

%%%%%%%%%%%%%%%%%%%%%%%%%%%%%%%%%%%%%%%%%%%%%%%%%%%%%%%
%% Preâmbulo do template UFC - Pós-Graduação         %%
%%                                                   %%
%% Adaptado a partir do template original UFC para   %%
%% trabalhos de disciplina de pós-graduação com foco %%
%% em Engenharia Elétrica / Eletrônica de Potência.  %%
%%%%%%%%%%%%%%%%%%%%%%%%%%%%%%%%%%%%%%%%%%%%%%%%%%%%%%%

% =====================================================
%   Codificação, idioma e tipografia
% =====================================================
\usepackage[utf8]{inputenc}                         % Acentuação direta
\usepackage[T1]{fontenc}                            % Codificação da fonte em 8 bits
\usepackage{microtype}                              % Melhorias finas de justificação
\usepackage{indentfirst}                            % Endenta o primeiro parágrafo de cada seção
\usepackage{mathptmx}                               % Fonte Times New Roman para texto e matemática

% =====================================================
%   Matemática
% =====================================================
\usepackage{amsfonts, amssymb, amsmath}             % Símbolos e ambientes matemáticos
\usepackage{amsgen}
\usepackage{mathtools}                              % Extensões de amsmath (matrizes, alinhamentos)

% Operadores em português (não há \sen, \tg, \cotg em LaTeX padrão)
\DeclareMathOperator{\sen}{sen}
\DeclareMathOperator{\tg}{tg}
\DeclareMathOperator{\cotg}{cotg}
\DeclareMathOperator{\arctg}{arctg}

% =====================================================
%   Figuras, tabelas e legendas
%   (memoir/abntex2 já fornecem [H] para floats e
%   ambiente subcaption via \newsubfloat — por isso
%   NÃO carregamos float nem subcaption aqui.)
% =====================================================
\usepackage{graphicx}                               % Inserir figuras
\graphicspath{{figuras/}{figuras/circuitos/}{figuras/formas-de-onda/}}
\usepackage{booktabs}                               % Tabelas profissionais (\toprule etc.)
\usepackage{multirow, array}                        % Múltiplas linhas e colunas em tabelas
\usepackage[font=singlespacing,format=plain,justification=justified,skip=0pt,singlelinecheck=false]{caption}

% =====================================================
%   Unidades SI e notação numérica (siunitx)
% =====================================================
\usepackage{siunitx}
\sisetup{
	output-decimal-marker = {,},
	per-mode               = symbol,
	inter-unit-product     = \ensuremath{{}\cdot{}},
	exponent-product       = \cdot,
	separate-uncertainty   = true,
	range-units            = single,
	range-phrase           = {\ \text{a}\ },
	list-units             = single,
	list-pair-separator    = {\ \text{e}\ },
	list-final-separator   = {\ \text{e}\ }
}

% =====================================================
%   Diagramas de circuitos (CircuiTikZ)
% =====================================================
\usepackage{tikz}
\usepackage[siunitx,european,straightvoltages,RPvoltages]{circuitikz}
\usetikzlibrary{calc,arrows.meta,positioning,patterns,decorations.markings}
\ctikzset{
	bipoles/length=1.0cm,
	diodes/scale=0.7,
	transistors/scale=0.8,
	resistors/scale=0.7,
	capacitors/scale=0.7,
	inductors/scale=0.7,
	sources/scale=0.8
}

% =====================================================
%   Gráficos vetoriais (pgfplots) — formas de onda,
%   curvas de FP e THD em função de Po, etc.
% =====================================================
\usepackage{pgfplots}
\pgfplotsset{
	compat=1.18,
	every axis/.append style={
		font=\small,
		line width=0.6pt,
		axis lines=left,
		grid=both,
		major grid style={line width=.2pt,draw=gray!30},
		minor grid style={line width=.1pt,draw=gray!15},
		tick align=outside
	},
	every axis plot/.append style={line width=1pt}
}
\usepgfplotslibrary{groupplots,fillbetween}

% =====================================================
%   Listagens de código e algoritmos
% =====================================================
\usepackage{verbatim}
\usepackage[svgnames,table]{xcolor}                 % Cores extras (Green, Blue, etc.)
\usepackage{listings}
\usepackage[portuguese,ruled,lined]{algorithm2e}
\usepackage{algorithmic}

% =====================================================
%   Estrutura, formatação e auxiliares
% =====================================================
\usepackage{tocloft}                                % Personaliza o Sumário
\usepackage{etoolbox}                               % Toolbox de comandos
\usepackage[nogroupskip,nonumberlist]{glossaries}   % Lista de abreviaturas e siglas
\usepackage[alf, abnt-emphasize=bf, recuo=0cm, abnt-etal-cite=2, abnt-etal-list=0, abnt-etal-text=it]{abntex2cite}
\usepackage{lipsum}                                 % Texto fictício (Lorem Ipsum)
\usepackage{appendix}                               % Geração de apêndices
\usepackage{paracol}                                % Parágrafos sem identação
\usepackage{pdfpages}                               % Inclui PDFs externos
\usepackage{lib/ufctex}                             % Normas UFC

% =====================================================
%   Hyperref — já é carregado pelo abntex2; apenas
%   complementamos a configuração para melhor SyncTeX
%   e bookmarks bem ordenados no leitor de PDF.
% =====================================================
\hypersetup{
	bookmarksopen        = true,
	bookmarksnumbered    = true,
	bookmarksopenlevel   = 1,
	unicode              = true,
	pdfencoding          = auto,
	pdfstartview         = FitH,
	breaklinks           = true,
	colorlinks           = false
}

% =====================================================
%   Glossário e índice
% =====================================================
\makeglossaries
\makeindex

\setlength{\mathindent}{0pt}

\trabalhoacademico{posgraduacao}
\ehqualificacao{nao}
\removerbordasdohyperlink{sim}
\cordohyperlink{nao}

\ies{Universidade Federal do Ceará}
\iessigla{UFC}
\centro{Centro de Tecnologia}
\departamento{Departamento de Engenharia Elétrica}
\programamestrado{Grupo de Processamento de Energia e Controle (GPEC)}
\disciplina{}          % vazio: a capa NAO imprime "Disciplina:" (e um relatorio)
\disciplinacodigo{}
\professor{}
\professorfeminino{nao}

\autor{João Rodrigo Arnaud da Cruz\\Adolfo Araújo}
\titulo{Ferramenta de Densidade de Potência para Conversores CC--CC Isolados: Relatório Técnico dos Modelos de Perdas, Magnéticos e Dissipadores, com Dedução das Equações}
\data{2026}
\local{Fortaleza}

\begin{document}
	\imprimircapa

	\begin{resumo}
		Este relatório técnico documenta, de forma completa e rastreável, a
		fundamentação teórica da \textbf{Ferramenta de Densidade de Potência}
		(GPEC/UFC), que estima a densidade de potência (\si{\watt\per\liter}) e o
		rendimento (\si{\percent}) de um conversor CC--CC isolado varrendo
		combinações de \textbf{semicondutores SiC} (MOSFET e diodo),
		\textbf{indutores}, \textbf{transformadores} e \textbf{dissipadores}.
		Para cada família apresentam-se o modelo de perdas e o procedimento de
		dimensionamento, com a \textbf{dedução de cada equação a partir das leis
		físicas} (Faraday, Ampère, Ohm térmico, Joule) e a \textbf{justificativa
		de cada constante}. Detalham-se a interpolação dos modelos de comutação
		PLECS, o modelo linear por partes do diodo, a analogia térmica do
		dissipador com seus dois casos-limite, o algoritmo iterativo de projeto
		do indutor com o \emph{roll-off} de permeabilidade, a lei de Faraday
		generalizada por uma constante de forma de onda~$k$ ajustável, as perdas
		de núcleo por Steinmetz e por rede neural (com normalização), o Steinmetz
		generalizado (GSE) e a combinação multiplicativa que produz o diagrama
		eficiência--densidade. Fecha-se com os parâmetros do protótipo de
		referência. O objetivo é permitir que a comunidade científica
		\textbf{use, verifique e aprimore} o algoritmo.
		\palavraschave{Densidade de potência. Conversor CC--CC isolado. SiC.
		Lei de Faraday. Steinmetz. Dissipador. Rendimento. Código aberto.}
	\end{resumo}

	\pdfbookmark[0]{\contentsname}{toc}
	\imprimirsumario

	\textual

%=====================================================================
\chapter{Introdução, metodologia e notação}
%=====================================================================
A \emph{densidade de potência} mede quanta potência um conversor processa por
unidade de volume; quanto maior, mais compacto e leve o conversor. A ferramenta
avalia um conversor CC--CC isolado composto por cinco famílias de componentes
--- dois conjuntos de MOSFETs, diodos, indutor de filtro e transformador --- mais
os dissipadores que dão suporte térmico aos semicondutores.

\section{Metodologia de cálculo}
O procedimento tem três etapas, aplicadas a cada família e depois ao conjunto:
\begin{enumerate}
	\item \textbf{Varredura.} Para cada componente do catálogo, e para cada ponto
	de operação (frequência, densidade de fluxo, densidade de corrente etc.),
	calculam-se as perdas e o volume, rejeitando os projetos que violam limites
	térmicos, de saturação ou de janela.
	\item \textbf{Filtragem.} O usuário seleciona subconjuntos (material,
	fabricante, encapsulamento, tipo de convecção) e a interface reduz o conjunto
	aos projetos de interesse.
	\item \textbf{Combinação.} Escolhido um projeto de cada família, somam-se
	perdas e volumes (Capítulo~\ref{cap:comb}) e calculam-se o rendimento e a
	densidade do conversor completo.
\end{enumerate}
Em cada família, o projeto ótimo é o par $(V,\,P_{\mathrm{perdas}})$ mais próximo
da \textbf{origem} (simultaneamente menor volume e menor perda).

\section{Figuras de mérito}
A \textbf{densidade de potência} é definida por
\begin{equation}
	\rho \;=\; \frac{P_{\mathrm{out}}}{V_{\mathrm{total}}}\qquad[\si{\watt\per\liter}].
	\label{eq:rho}
\end{equation}
O rendimento sai da própria definição, $\eta=P_{\mathrm{out}}/P_{\mathrm{in}}$.
Pela conservação de energia, toda a potência de entrada ou vira saída ou vira
perda: $P_{\mathrm{in}}=P_{\mathrm{out}}+P_{\mathrm{perdas}}$. Tratando
$P_{\mathrm{out}}$ como a potência de referência do projeto,
\begin{equation}
	\eta=\frac{P_{\mathrm{out}}}{P_{\mathrm{in}}}
	    =\frac{P_{\mathrm{out}}}{P_{\mathrm{out}}+P_{\mathrm{perdas}}}
	    =\frac{P_{\mathrm{out}}-P_{\mathrm{perdas}}}{P_{\mathrm{out}}}\times100\%,
	\label{eq:eta}
\end{equation}
sendo a última forma a adotada na ferramenta (idêntica à do aplicativo original).
Como $\eta$ contém $P_{\mathrm{out}}$, ela depende da potência nominal; uma
análise por perdas absolutas (\si{\watt}) ou por densidade de perdas
(\si{\watt\per\liter}) independe dela.

\section{Notação}
$V_d$: tensão CC aplicada ao componente; $f_s$: frequência de comutação;
$T_j,\,T_{amb}$: temperaturas de junção e ambiente; $I_{\mathrm{ef}},\,I_{\mathrm{med}}$:
correntes eficaz e média; $A_e$: área efetiva do núcleo; $\ell_e$: comprimento do
caminho magnético; $V_e$: volume efetivo; $B_m$: densidade de fluxo de pico; $N$:
número de espiras; $\mu_0=4\pi\times10^{-7}\,\si{\henry\per\meter}$.

%=====================================================================
\chapter{Perdas nos semicondutores}
%=====================================================================
Os semicondutores do catálogo são de \textbf{carbeto de silício} (SiC), de
Wolfspeed e GeneSiC. Perdas em SiC dividem-se em comutação e condução.

\section{MOSFET: comutação}
Do ponto de vista físico, durante uma transição a tensão e a corrente do dispositivo
se sobrepõem por um curto intervalo; a energia dissipada é a integral da potência
instantânea sobre a transição:
\begin{equation}
	E_{sw}=\int_{\text{transição}} v_{ds}(t)\,i_d(t)\,dt\quad[\si{\joule}].
	\label{eq:esw_int}
\end{equation}
Como essa integral depende de muitos fatores não lineares, adota-se a solução da
indústria: usam-se \textbf{modelos térmicos PLECS} (tabelas dos
\emph{datasheets}), interpolados como função da corrente comutada~$I$, da tensão
de bloqueio~$V_{ds}$, da temperatura de junção~$T_j$ e do resistor de porta~$R_g$:
\begin{equation}
	E_{on}=f_{on}(I_{on},V_{ds},T_j,R_g),\qquad
	E_{off}=f_{off}(I_{off},V_{ds},T_j,R_g).
	\label{eq:eon}
\end{equation}
\textit{Dependências.} A energia cresce aproximadamente com $V_{ds}\,I$ (área do
retângulo tensão--corrente) e com $R_g$ (que retarda $dv/dt$ e $di/dt$, alongando
a sobreposição); cresce ainda com $T_j$ (mobilidade e cauda de corrente).
\textit{Dedução da potência.} Se cada evento dissipa $E$ e ocorre uma vez por
período $T_s=1/f_s$, a potência média é a energia por evento vezes a taxa de
eventos:
\begin{equation}
	P_{on}=\frac{E_{on}}{T_s}=f_s\,E_{on},\qquad P_{off}=f_s\,E_{off}.
	\label{eq:pon}
\end{equation}
Sob \emph{Zero-Voltage Switching} (ZVS) a entrada comuta com $v_{ds}=0$, anulando
a integral~\eqref{eq:esw_int} de entrada: $P_{on}=0$.

\section{MOSFET: condução}
\textit{Dedução.} No estado ligado o canal SiC comporta-se como resistência,
$v=R_{ds,on}\,i$. A potência instantânea é $p(t)=v\,i=R_{ds,on}\,i^2$ e a média em
um período vale
\begin{equation}
	P_{cond}=\frac{1}{T_s}\!\int_0^{T_s}\!R_{ds,on}\,i^2\,dt
	=R_{ds,on}\,\langle i^2\rangle=R_{ds,on}\,I_{\mathrm{ef}}^{2},
	\label{eq:pcondm}
\end{equation}
pois, por definição, $I_{\mathrm{ef}}^2=\langle i^2\rangle$. A resistência
equivalente é estimada pela média de $V_{drop}/i$ nos instantes de ligar e
desligar --- como $V_{drop}=i\,R_{ds}$, a razão recupera $R_{ds}$:
\begin{equation}
	R_{ds,on}=\tfrac12\!\left[\frac{V_{drop}(I_{on},V_{gs},T_j)}{I_{on}}
	+\frac{V_{drop}(I_{off},V_{gs},T_j)}{I_{off}}\right].
	\label{eq:rds}
\end{equation}
A aproximação é válida porque, no SiC, a característica $v(i)$ do estado ligado é
quase linear. A perda total do MOSFET é $P_{\mathrm{total}}=P_{on}+P_{off}+P_{cond}$.

\section{Diodo SiC Schottky}
\textit{Modelo.} A curva direta de um diodo, exponencial em pequena escala, é
aproximada por um modelo linear por partes: uma fonte de limiar $V_{TO}$ em série
com uma resistência dinâmica $R_d$, ambas quadráticas na temperatura:
\begin{align}
	R_d(T_j)  &= R_{d,a}\,T_j^{2}+R_{d,b}\,T_j+R_{d,c},&
	V_{TO}(T_j)&= V_{TO,a}\,T_j^{2}+V_{TO,b}\,T_j+V_{TO,c}.
	\label{eq:rdvto}
\end{align}
\textit{Dedução da perda.} Com $v(i)=V_{TO}+R_d\,i$, a potência instantânea é
$p=v\,i=V_{TO}\,i+R_d\,i^2$; tomando a média e usando
$\langle i\rangle=I_{\mathrm{med}}$ e $\langle i^2\rangle=I_{\mathrm{ef}}^2$:
\begin{equation}
	P_{cond}=V_{TO}\,I_{\mathrm{med}}+R_d\,I_{\mathrm{ef}}^{2}.
	\label{eq:pcondd}
\end{equation}
Note que a parcela de limiar pondera a corrente \emph{média} e a resistiva a
\emph{eficaz}, refletindo as respectivas leis $p=V i$ e $p=R i^2$. O diodo
Schottky de SiC \textbf{não} tem carga de recuperação reversa relevante
($Q_{rr}\approx0$), logo não há parcela de comutação. Se o encapsulamento contém
dois \emph{dies} em meia-ponte, a corrente se divide e
$P_{cond}=2[R_d(I_{\mathrm{ef}}/2)^2+V_{TO}(I_{\mathrm{med}}/2)]$.

%=====================================================================
\chapter{Dimensionamento do dissipador}
%=====================================================================
Por analogia com o circuito elétrico: em regime permanente, a condução de calor obedece a
uma lei análoga à de Ohm: o fluxo de calor $P$ (\si{\watt}) faz o papel da
corrente, a diferença de temperatura $\Delta T$ o da tensão, e a resistência
térmica $R_\theta$ (\si{\degreeCelsius\per\watt}) o da resistência elétrica, com
$\Delta T=P\,R_\theta$. As três interfaces em série somam-se:
\begin{equation}
	T_j-T_{amb}=P_{\mathrm{loss}}\,(R_{jc}+R_{chs}+R_{ha})
	\;\Rightarrow\;
	T_j = T_{amb}+P_{\mathrm{loss}}\,(R_{jc}+R_{chs}+R_{ha}).
	\label{eq:chain}
\end{equation}
Usam-se $R_{jc}=\SI{1}{\degreeCelsius\per\watt}$ e
$R_{chs}=\SI{0.5}{\degreeCelsius\per\watt}$ como aproximações genéricas de
TO-247/SOT-227 (poderiam vir do \emph{datasheet}). Isolando a máxima temperatura
admissível do dissipador pelo lado do dispositivo:
\begin{equation}
	T_{hs,\mathrm{dev}} = T_j - P_{\mathrm{loss}}\,(R_{jc}+R_{chs}).
	\label{eq:ths}
\end{equation}
\textit{Dedução da resistência necessária.} Com $Q$ dispositivos por dissipador,
a potência total nele é $Q\,P_{\mathrm{loss}}$; a queda dissipador--ambiente é
$Q\,P_{\mathrm{loss}}\,R_{ha}$. Igualando à elevação disponível surgem dois
\textbf{casos-limite}, e adota-se o mais restritivo:
\begin{equation}
	R_{ha,\mathrm{req}}=\min\!\left\{
	\frac{T_{hs,\mathrm{dev}}-T_{amb}}{Q\,P_{\mathrm{loss}}}\ \text{(limite da junção)},\;
	\frac{\Delta T_{\mathrm{adm}}}{Q\,P_{\mathrm{loss}}}\ \text{(limite do dissipador)}
	\right\},
	\label{eq:rhareq}
\end{equation}
onde $\Delta T_{\mathrm{adm}}=T_{hs,\max}-T_{amb}$. \textit{Comprimento.} A
relação comprimento~$\times$~resistência dos perfis não é analítica; ajusta-se
empiricamente a curva do catálogo HSDissipadores, isolando o comprimento, com
$\mathrm{coef}=R_{ha,\mathrm{req}}/R_{ha}$:
\begin{equation}
	L = \frac{64{,}1030}{\mathrm{coef}-0{,}4258}+14{,}754\quad[\si{\milli\meter}].
	\label{eq:Lfit}
\end{equation}
As três constantes são os parâmetros dessa regressão (mantidos idênticos ao
projeto original). O volume é o produto das três dimensões:
\begin{equation}
	V_{hs}=\frac{L}{10}\cdot\frac{H}{10}\cdot\frac{W}{10}\quad[\si{\cubic\centi\meter}].
	\label{eq:vhs}
\end{equation}
É esse volume que entra na densidade, pois o do semicondutor é desprezível.

%=====================================================================
\chapter{Projeto do indutor}
%=====================================================================
Indutor de filtro toroidal com núcleo de pó. O projeto é iterativo porque a
permeabilidade cai com a polarização CC.

\section{Indutância e corrente eficaz}
\textit{Dedução de $L$.} De $v=L\,di/dt$, integrando no meio-período em que a
tensão média $V_d$ é aplicada, a excursão de corrente é $\Delta I_L$. Para razão
cíclica $D=0{,}5$ (pior caso do \emph{ripple}), o resultado padrão do conversor
comutado é $\Delta I_L=V_d\,D(1-D)/(Lf_s)=V_d/(4Lf_s)$; isolando $L$:
\begin{equation}
	L=\frac{V_d}{4\,f_s\,\Delta I_L}.
	\label{eq:Lind}
\end{equation}
\textit{Dedução de $I_{L,\mathrm{ef}}$.} A corrente é $i(t)=I_{L,\mathrm{med}}+i_{\sim}(t)$,
com $i_{\sim}$ triangular de amplitude de pico $\Delta I_L/2$ e média nula. Como a
componente CC e a alternada são ortogonais, $\langle i^2\rangle=I_{L,\mathrm{med}}^2+\langle i_{\sim}^2\rangle$.
O valor quadrático médio de uma triangular simétrica de pico $\hat a$ é
$\langle i_\sim^2\rangle=\hat a^2/3$; com $\hat a=\Delta I_L/2$ resulta
$\langle i_\sim^2\rangle=\Delta I_L^2/12$, logo
\begin{equation}
	I_{L,\mathrm{ef}}=\sqrt{I_{L,\mathrm{med}}^2+\tfrac{\Delta I_L^2}{12}}
	=I_{L,\mathrm{med}}\sqrt{1+\tfrac{1}{12}\!\left(\tfrac{\Delta I_L}{I_{L,\mathrm{med}}}\right)^{2}}.
	\label{eq:ILef}
\end{equation}

\section{Espiras e \emph{roll-off} de permeabilidade}
\textit{Dedução de $A_{L0}$.} Um circuito magnético de comprimento $\ell_e$, área
$A_e$ e permeabilidade $\mu_0\mu_{\mathrm{rel}}$ tem relutância
$\mathcal{R}=\ell_e/(\mu_0\mu_{\mathrm{rel}}A_e)$. A indutância de $N$ espiras é
$L=N^2/\mathcal{R}$; a indutância por espira ao quadrado é
\begin{equation}
	A_{L0}=\frac{L}{N^2}=\mu_{\mathrm{rel}}\,\mu_0\,\frac{A_e}{\ell_e},\qquad
	N=\bigl\lceil\sqrt{L/A_L}\bigr\rceil.
	\label{eq:al0}
\end{equation}
\textit{Dedução do campo.} Pela lei de Ampère, $\oint\!H\,dl=N I\Rightarrow H\ell_e=NI$,
logo $H=NI/\ell_e$; convertendo à unidade prática (oersted, $\ell_e$ em mm) pela
relação CGS,
\begin{equation}
	H=4\pi\,\frac{N}{\ell_e}\,I_{L,\mathrm{med}}\quad[\mathrm{Oe}].
	\label{eq:H}
\end{equation}
\textit{Algoritmo iterativo.} A permeabilidade efetiva decresce com $H$ segundo o
ajuste (empírico, por material e permeabilidade) do \emph{roll-off}:
\begin{equation}
	\mu\%(H)=\frac{1}{a+b\,H^{c}}\times100\ \in[0,100].
	\label{eq:rolloff}
\end{equation}
Como $H$ depende de $N$ e $N$ depende de $A_L=A_{L0}\,\mu\%/100$, resolve-se por
ponto fixo:
\begin{algorithmic}[1]
	\STATE $A_L\leftarrow A_{L0}$;\quad $N\leftarrow\lceil\sqrt{L/A_L}\rceil$
	\REPEAT
		\STATE $H\leftarrow 4\pi\,(N/\ell_e)\,I_{L,\mathrm{med}}$
		\STATE $\mu\%\leftarrow$ \eqref{eq:rolloff};\quad se $\mu\%<30$: \textbf{rejeita} (saturado)
		\STATE $A_L\leftarrow A_{L0}\,\mu\%/100$;\quad $N\leftarrow\lceil\sqrt{L/A_L}\rceil$
	\UNTIL{$N$ estabiliza}
\end{algorithmic}

\section{Densidade de fluxo e perdas}
\textit{Dedução de $B_{pk}$.} O fluxo enlaçado é $\lambda=N\Phi=L\,I$ e o fluxo
$\Phi=B\,A_e$; logo $B=\Phi/A_e=\lambda/(N A_e)$:
\begin{equation}
	B_{pk}=\frac{L\,I}{N\,A_e},\qquad I=I_{L,\mathrm{med}}\pm\tfrac{\Delta I_L}{2}.
	\label{eq:bpk}
\end{equation}
Esta identidade usa $L$ e o inteiro $N$ já obtidos, sem recorrer à permeabilidade
(que divergiria). O aplicativo legado usava o comprimento médio de espira no lugar
de $\ell_e$ e omitia $N$, subestimando $B$ (e as perdas) em \numrange{4}{5}~vezes.

\textit{Perda de núcleo (Steinmetz).} A perda volumétrica de um material
magnético segue empiricamente uma lei de potência na frequência e no fluxo (lei
de Steinmetz), com expoentes ajustados por material:
\begin{equation}
	P_v = P_a\,B_{pk}^{\,P_b}\,\!\left(\frac{f_s}{1000}\right)^{\!P_c}\ [\si{\milli\watt\per\cubic\centi\meter}],
	\quad P_{\mathrm{core}}=P_v\cdot10^{-3}\,V_{\mathrm{nucleo}}\ [\si{\watt}].
	\label{eq:pcore}
\end{equation}
\textit{Perda no cobre.} A resistência é $R=\rho_{cu}\,\ell_{fio}/A_{fio}$ com
$\ell_{fio}=N\,\ell_{\mathrm{esp}}$ e $N_{\mathrm{paral}}$ fios em paralelo; como
$P=R I^2$ e a resistividade do cobre cresce ($R(T)=R_{20}[1+0{,}00393(T-20)]$):
\begin{equation}
	P_{cu}=\frac{N\,\ell_{\mathrm{esp}}\,R_{awg}(T)}{N_{\mathrm{paral}}}\,I_{L,\mathrm{ef}}^{2}.
	\label{eq:pcu}
\end{equation}
\textit{Elevação de temperatura.} A convecção natural obedece à relação empírica
$\Delta T=(10^{5}P/S_a)^{0{,}833}$, com $S_a$ a área de superfície. O expoente
$0{,}833\approx5/6$ é típico de placas em convecção natural. \textit{Utilização da
janela.} $K_u=A_{\mathrm{ocup}}/A_w$ com
$A_{\mathrm{ocup}}=\tfrac{\pi}{4}D_{fio}^2\,N\,N_{\mathrm{paral}}$, exigindo
$K_{u,\min}\le K_u\le K_{u,\max}$.

%=====================================================================
\chapter{Projeto do transformador}
%=====================================================================
Transformador de alta frequência em ferrites Magnetics (Steinmetz), TDK (rede
neural) ou Magmattec.

\section{Lei de Faraday: espiras para um dado \texorpdfstring{$B_{\max}$}{Bmax}}
\textit{Dedução completa.} Pela lei de Faraday, a tensão no enrolamento é
$v(t)=N\,d\Phi/dt=N A_e\,dB/dt$. Integrando no meio-período $[\,0,\,1/(2f)\,]$ em
que a tensão $V$ é aplicada, o fluxo excursiona de $-B_m$ a $+B_m$:
\begin{equation}
	\int_0^{1/(2f)}\! v\,dt = N A_e\!\int_{-B_m}^{+B_m}\! dB = N A_e\,(2B_m).
	\label{eq:vs}
\end{equation}
Para uma onda com valor característico $V$ aplicado durante esse meio-período, o
lado esquerdo vale $V/(2f)$ multiplicado por um fator de forma; reagrupando os
fatores numéricos numa única constante $k$ chega-se à forma canônica
\begin{equation}
	N=\frac{V}{k\,f\,B\,A_e},\qquad
	k=\begin{cases}4 & \text{onda quadrada}\\ 4{,}44=\sqrt2\,\pi & \text{senoidal}\\ 9 & \text{inversor 3 níveis}\end{cases}
	\label{eq:Ngeral}
\end{equation}
\textit{Justificativa dos valores.} Para onda quadrada simétrica,
$V=N A_e(2B_m)(2f)=4NA_eB_mf\Rightarrow k=4$. Para senoide, a relação entre valor
eficaz e de pico introduz o fator $\sqrt2\,\pi\approx4{,}44$. O valor $k=9$ é
específico da forma de onda do inversor trifásico de 3 níveis deste projeto
(herdado do aplicativo e do artigo de referência). Com $A_e$ em
\si{\square\centi\meter} ($A_{e,\mathrm{m^2}}=A_{e,\mathrm{cm^2}}\times10^{-4}$), o
$10^{-4}$ passa ao numerador como $10^{4}$:
\begin{equation}
	\boxed{\;N_p=\operatorname{round}\!\left(\frac{V_d\cdot10^{4}}{k\,B_m\,f_s\,A_{e,\mathrm{cm^2}}}\right)\;},
	\qquad N_s=\lceil N_p\,n\rceil.
	\label{eq:Np}
\end{equation}
Na ferramenta, $k$ é um \textbf{campo de entrada} (``Faraday Constant''), padrão
$9$: o usuário informa a constante conforme a forma de onda. Menos espiras
implicam $B$ maior (risco de saturação); mais espiras reduzem $B$ mas ocupam a
janela --- por isso a varredura testa vários $B_m$.

\section{Perdas de núcleo: Steinmetz e rede neural}
Para ferrites Magnetics/Magmattec, usa-se Steinmetz com correção térmica
quadrática:
\begin{equation}
	nP_v=\frac{k_s\,f_s^{\alpha}\,B_m^{\beta}\,(c_1-c_2 T+c_3 T^{2})}{1000},\qquad
	P_{\mathrm{core}}=V_e\,\frac{nP_v}{1000}.
	\label{eq:steintraf}
\end{equation}
A ordem correta dos coeficientes é $k_s,\alpha,\beta$; trocá-los produziria
$f_s^{\alpha}$ com o expoente de $k_s$ e perdas irreais.
\textit{Rede neural (TDK N27--N97).} Para esses ferrites, $nP_v$ vem de uma rede
\emph{feed-forward} pré-treinada. As entradas $(f_s,B_m,T)$ são normalizadas ao
intervalo de treino, $\hat x=(x-x_{\min})/(x_{\max}-x_{\min})$; a camada oculta
aplica $\tanh(W\hat x+b)$; a saída é desnormalizada de volta a
\si{\milli\watt\per\cubic\centi\meter}. Pesos $W$ e \emph{bias} $b$ vêm de um
arquivo \texttt{.mat} pré-treinado.

\section{Steinmetz generalizado (GSE)}
\textit{Motivação.} Steinmetz clássico vale para excitação senoidal; para a forma
de onda real, integra-se a taxa instantânea de variação do fluxo (GSE):
\begin{equation}
	P_{GSE}=\frac{1}{T}\!\int_0^{T} k_i\,\Bigl|\tfrac{dB}{dt}\Bigr|^{\alpha}\,|B|^{\beta-\alpha}\,dt,
	\quad
	k_i=\frac{k_s}{(2\pi)^{\alpha-1}\!\displaystyle\int_0^{2\pi}\!|\cos\theta|^{\alpha}|\sin\theta|^{\beta-\alpha}d\theta}.
	\label{eq:gse}
\end{equation}
O denominador de $k_i$ garante que, para excitação senoidal, o GSE reproduza o
Steinmetz clássico. É desativado na varredura (custo das integrais) e retorna
indefinido se $\beta-\alpha\le-0{,}99$ (integral divergente).

\section{Janela, cobre e temperatura}
$K_u=A_{\mathrm{nec}}/A_w$ com $A_w=A_p/A_e$, restrito a
$[K_{u,\min},K_{u,\max}]$. A perda no cobre soma primário e secundário e a
elevação segue a mesma lei de convecção:
\begin{equation}
	P_{cond}=R_p\,I_{p,\mathrm{ef}}^{2}+R_s\,I_{s,\mathrm{ef}}^{2},\qquad
	\Delta T=\!\left(\frac{1000\,(P_{\mathrm{core}}+P_{cond})}{S_a}\right)^{\!0{,}833}.
	\label{eq:dttraf}
\end{equation}
O volume total é $V_{\mathrm{total}}=V_e+V_{cu}$; rejeita-se se $\Delta T$, a
perda ou o volume excederem os limites.

%=====================================================================
\chapter{Combinação e densidade do conjunto}
\label{cap:comb}
%=====================================================================
A aba final combina uma escolha de cada família. \textit{Modelo de unidade.} A
unidade de um semicondutor é \textbf{um dissipador com seus $Q$ \emph{dies}}: a
perda da unidade é a do conjunto ($Q$ vezes a de um \emph{die}, pois todos os
$Q$ aquecem o mesmo dissipador) e o volume é o de um dissipador; o multiplicador
$m_i$ é o número de braços do conversor. Somando sobre as famílias:
\begin{equation}
	P_{\mathrm{perdas}}=\sum_i m_i\,P_{\mathrm{unidade},i},\qquad
	V_{\mathrm{total}}=\sum_i m_i\,V_{\mathrm{unidade},i}.
	\label{eq:combina}
\end{equation}
\textit{Dedução das densidades.} Dividindo a potência (de saída ou de perdas)
pelo volume total em litros ($1\,\si{\liter}=10^{3}\,\si{\cubic\centi\meter}$):
\begin{equation}
	\rho_{P}=\frac{P_{\mathrm{out}}}{V_{\mathrm{total}}\cdot10^{-3}},\qquad
	\rho_{L}=\frac{P_{\mathrm{perdas}}}{V_{\mathrm{total}}\cdot10^{-3}}\quad[\si{\watt\per\liter}].
	\label{eq:densi}
\end{equation}
\textit{Complexidade e amostragem.} A combinação de cinco famílias é
multiplicativa; se cada uma tem $n_i$ projetos, há $\prod_i n_i$ conjuntos ---
número astronômico. Amostra-se, portanto, um subconjunto representativo por
família (espalhado pela faixa de perdas, para preservar a fronteira de
compromisso), ou, sob confirmação do usuário, combinam-se todos até um limite de
memória. A vista padrão grafa Eficiência~$\times$~$\rho_L$ (reproduz o aplicativo
original) e a alternativa Eficiência~$\times$~$\rho_P$.

%=====================================================================
\chapter{Parâmetros do protótipo de referência}
%=====================================================================
Os valores padrão da interface derivam do conversor CC--CC isolado trifásico do
artigo de referência (SOBRAEP/OJPE), na Tabela~\ref{tab:proto}.

\begin{table}[htb]
	\centering
	\caption{Especificações do conversor de referência.}
	\label{tab:proto}
	\begin{tabular}{@{}lll@{}}
		\toprule
		Grandeza & Símbolo & Valor \\
		\midrule
		Tensão de entrada        & $V_{in}$ & \SI{600}{\volt} \\
		Tensão de saída          & $V_{o}$  & \SI{920}{\volt} \\
		Potência de saída        & $P_{o}$  & \SI{10}{\kilo\watt} \\
		Corrente de saída        & $I_{o}$  & \SI{14.28}{\ampere} \\
		Frequência de comutação  & $f_s$    & \SI{50}{\kilo\hertz} \\
		Indutor de comutação     & $L_d$    & \SI{5}{\micro\henry} \\
		Relação de espiras       & $n$      & $0{,}8333$ \\
		Indutor de filtro        & $L_f$    & \SI{1}{\milli\henry} \\
		Capacitor de saída       & $C_o$    & \SI{100}{\micro\farad} \\
		\bottomrule
	\end{tabular}
\end{table}

No protótipo otimizado, o indutor de filtro usa núcleo Magnetics MPP de
permeabilidade $\mu=147$ com 52 espiras, e o transformador usa núcleo Magmattec
MMT139 com $N_p=35$ e $N_s=42$. A tensão sobre o indutor de filtro é de cerca de
\SI{300}{\volt} (não o barramento), e a corrente eficaz do primário é de
aproximadamente \SI{2.8}{\ampere} --- fatos que justificam os valores padrão e
evitam rejeições espúrias por saturação ou por janela cheia.

%=====================================================================
\chapter{Conclusão}
%=====================================================================
Apresentou-se a fundamentação teórica completa da ferramenta, com a
\textbf{dedução de cada equação a partir das leis físicas} e a
\textbf{justificativa de cada constante}: as energias e potências de comutação e
condução dos semicondutores SiC; a analogia térmica e os dois casos-limite do
dissipador; o projeto iterativo do indutor com a identidade exata de $B_{pk}$, o
algoritmo de \emph{roll-off} e Steinmetz por material; o projeto do transformador
com a \textbf{lei de Faraday generalizada} pela constante $k$ ajustável,
incluindo Steinmetz, rede neural e GSE; e a combinação multiplicativa que produz
o diagrama eficiência--densidade. Todos os resultados numéricos da interface são,
assim, rastreáveis à teoria de eletrônica de potência, de modo que a comunidade
científica possa reproduzi-los, verificá-los e aprimorá-los.

\end{document}
